%% ==========================================================================
%%  4  Unit subsidies for objective signed dichotomous valuations
%%
%%  Content of Section 4 of the arXiv version (arxiv/sections/objective_mixed.tex)
%%  and its algorithm (arXiv Appendix C), in the house style. The proof of
%%  Lemma~\ref{lem:goods-extension} cites the algorithm of Barman et al.
%%  directly (2026-10-08). The step-by-step check that their argument survives
%%  chores in the bundles is Appendix~\ref{app:goods-extension}.
%% ==========================================================================
\section{Unit Subsidies for Objective Signed Dichotomous Valuations}
\label{sec:objective-mixed}

We now consider mixed instances in which every item is commonly classified as a
good or a chore. Such common-sign settings are referred to as \emph{objective}
goods and chores or \emph{objective valuations} in the fair-division
literature~\cite{aziz2022fair,bilo2026approximately}. We additionally impose
signed binary marginals.

\begin{definition}[Objective signed dichotomous valuations]
\label{def:objective-signed}
A valuation profile $\set{v_i}_{i \in \agents}$ is \emph{objective signed
dichotomous} if the item set admits a common partition
$\items = G \mathbin{\dot\cup} C$ such that, for every agent $i \in \agents$ and
every $S \subseteq \items$,
\[
  v_i(S\cup\{g\})-v_i(S)\in\{0,1\}
  \qquad
  \text{for every }g\in G\setminus S,
\]
and
\[
  v_i(S\cup\{c\})-v_i(S)\in\{0,-1\}
  \qquad
  \text{for every }c\in C\setminus S.
\]
\end{definition}

Thus, every item is commonly regarded as either a good or a chore, although its
marginal value may depend on the bundle to which it is added. In particular, the
valuations need not be additive, submodular, or subadditive. The class is a
subclass of doubly monotone valuations in which the classification of items into
goods and chores is common across agents. Since $v_i(\emptyset) = 0$ and every
marginal is an integer, the telescoping argument of
Observation~\ref{obs:integrality} shows that every such valuation is integer
valued.

The plan is to allocate the chores by Theorem~\ref{thm:m-main} and then the goods
by the algorithm of Barman et al.~\cite{BKNS22}. That algorithm is incremental:
it inserts the goods one at a time, and its only invariant is that the current
solution is envy-free with subsidies in $\{0,1\}^n$. So it can take over from any
solution satisfying the invariant, provided its one-good step remains valid when
the bundles already contain chores. The following lemma records that step.

\begin{lemma}
\label{lem:goods-extension}
Let $M'$ be a set of allocated items and let $A=(A_1,\ldots,A_n)$ be a complete
allocation of $M'$. Suppose that the valuations are integer-valued and that
there exists $p\in\{0,1\}^n$ such that $(A,p)$ is envy-free. Let $g\notin M'$
satisfy
\[
  v_i(S\cup\{g\})-v_i(S)\in\{0,1\}
\]
for every agent $i$ and every $S\subseteq M'$. Then there exist a complete
allocation $A'$ of $M'\cup\{g\}$ and $p'\in\{0,1\}^n$ such that $(A',p')$ is
envy-free. Moreover, $(A',p')$ can be computed in polynomial time given
value-oracle access.
\end{lemma}

\begin{proof}
This is one iteration of the algorithm of Barman et al.~\cite{BKNS22}. If
$(A,p)$ is extendable with $g$, that is, if some permutation of the bundles keeps
the solution envy-free and gives a maximally subsidised agent a bundle on which
$g$ has marginal value one, the subroutine \textsc{Extend} finds such a
permutation and agent, and that agent receives $g$ (Lemmas~7 and~8
of~\cite{BKNS22}). Otherwise the subroutine \textsc{FindSink} finds a recipient
for $g$ (Lemmas~9 to~11 of~\cite{BKNS22}). In both cases the pointwise-minimal
subsidy vector of the new allocation lies in $\{0,1\}^n$, and the step runs in
polynomial time.

These lemmas are stated for instances in which every item is a good, whereas the
bundles of $A$ may contain chores. Their proofs, however, use that hypothesis
only through two facts: the valuations are integer valued, and the good being
inserted has marginal value in $\{0,1\}$ for every agent on every bundle.
Monotonicity and dichotomy are invoked only for that good, never for the items
already allocated. Both facts hold here, so the step applies unchanged.
Appendix~\ref{app:goods-extension} restates the argument and checks each step
under these two facts.
\end{proof}

We can now obtain the unit-subsidy guarantee for objective signed dichotomous
valuations.

\begin{theorem}
\label{thm:objective-mixed}
For every instance with objective signed dichotomous valuations, there exist a
complete allocation $\alloc$ and a subsidy vector $p\in\{0,1\}^n$ such that
$(\alloc,p)$ is envy-free and
\[
  \sum_{i\in N} p_i\leq n-1.
\]
Moreover, such an allocation and subsidy vector can be computed in polynomial
time in the value-oracle model.
\end{theorem}

\begin{proof}
Restrict the instance to the chores $C$. By Theorem~\ref{thm:m-main}, we can
compute an envy-free solution $(\alloc^C,p^C)$ with $p^C\in\{0,1\}^n$. Starting
from this solution, insert the goods in $G$ one at a time using
Lemma~\ref{lem:goods-extension}. After all insertions, we obtain a complete
envy-free solution $(\alloc,p)$ with $p\in\{0,1\}^n$. If every component of $p$
equals one, subtract one from every subsidy. Otherwise some component is already
zero. Thus, after normalization, $\sum_i p_i\le n-1$. Polynomial-time
computability follows from Theorem~\ref{thm:m-main} and
Lemma~\ref{lem:goods-extension}.
\end{proof}

The resulting procedure is stated explicitly as Algorithm~\ref{alg:objective-mixed}.
For brevity, let \textsc{ChoreAllocate} denote the polynomial-time procedure of
\Cref{sec:main} (Algorithm~\ref{alg:main}) for negative dichotomous valuations.
Given a negative dichotomous chore instance, it returns a complete allocation
$\alloc$ and $p\in\{0,1\}^n$ such that $(\alloc,p)$ is envy-free and
$\sum_i p_i\le n-1$.

\begin{algorithm}[H]
\caption{Unit subsidies for objective signed dichotomous valuations}
\label{alg:objective-mixed}
\begin{algorithmic}[1]
\Require Agents $N=[n]$, items $M=G\mathbin{\dot\cup}C$, and value-oracle access
  to objective signed dichotomous valuations $\{v_i\}_{i\in N}$.
\Ensure A complete allocation $\alloc$ and $p\in\{0,1\}^n$ such that
  $(\alloc,p)$ is envy-free and $\sum_{i\in N}p_i\le n-1$.

\State Restrict each $v_i$ to subsets of $C$.
\State $(\alloc,p)\gets\textsc{ChoreAllocate}(N,C,\{v_i|_C\}_{i\in N})$.
\State Replace $p$ by the pointwise-minimal subsidy vector supporting $\alloc$.

\For{each $g\in G$}
    \State Apply the one-good extension procedure of
      Lemma~\ref{lem:goods-extension} to $(\alloc,p)$ and $g$, obtaining an
      allocation $\alloc'$ of the currently allocated items together with $g$
      and a subsidy vector $p'\in\{0,1\}^n$ such that $(\alloc',p')$ is
      envy-free.
    \State $\alloc\gets\alloc'$.
    \State Replace $p'$ by the pointwise-minimal subsidy vector supporting
      $\alloc$, and set $p\gets p'$.
\EndFor

\State \Return $(\alloc,p)$.
\end{algorithmic}
\end{algorithm}

The one-good extension in Line~5 is implemented by the
\textsc{Extend}/\textsc{FindSink} procedure of Barman et al.~\cite{BKNS22}, which
remains valid in the present setting by Lemma~\ref{lem:goods-extension}. After
every iteration, we take the pointwise-minimal supporting subsidy vector. Since
the valuations are integer-valued, this vector is integral, and
Lemma~\ref{lem:goods-extension} implies that all its components are at most one.
Moreover, pointwise minimality implies that at least one component is zero.
Hence, throughout the algorithm
\[
  p\in\{0,1\}^n
  \qquad\text{and}\qquad
  \sum_{i\in N}p_i\le n-1.
\]

The total-subsidy bound in Theorem~\ref{thm:objective-mixed} is tight. Indeed,
negative dichotomous chore instances form a subclass of objective signed
dichotomous valuations, and the identical binary additive instance with $n$
agents and $n-1$ chores of Example~\ref{ex:lowerbound} already requires a total
subsidy of at least $n-1$.

Theorem~\ref{thm:objective-mixed} concerns the subsidy guarantee only. Our
construction does not guarantee that the underlying allocation is EF1, since the
subsequent goods-allocation procedure of Barman et al.~\cite{BKNS22} does not in
general preserve such a guarantee. Pareto optimality also cannot be imposed in
general, since Proposition~\ref{prop:po-impossibility} applies to the pure-chore
subclass. Under additivity, however, Pareto optimality can be recovered together
with the same subsidy bound (Proposition~\ref{prop:additive-objective-po}).
